\section{State of the Art}
\label{sec:state_of_art}
The methodological contributions of this work touch four distinct bodies of literature. Because this master's thesis proposes a continuous bilevel formulation of an airline hub-and-spoke design problem in which sparsity-inducing penalties replace integer activation variables and entropy regularization replaces the explicit logit constraint, the review is organized around those building blocks. We first cover the classical hub-and-spoke network design literature, with an emphasis on airline applications. We then review entropy regularization and logit demand models, which provide the behavioral foundation of the lower-level continuous formulation. We subsequently examine continuous and sparsity-based approximations of fixed-cost models, which underpin the lower-level problem. We close with a dedicated review of bilevel optimization and penalty methods, since the upper-level layer of the proposed model relies heavily on that machinery.

\subsection{Hub-and-Spoke Airline Network Design}
Hub-and-spoke network design has traditionally been modeled as a discrete location-and-routing problem in which the main decisions are the selection of hub facilities, the assignment of spoke nodes to those hubs, and the activation of direct or inter-hub services. The foundational formulations cast these decisions as mixed-integer optimization problems in which the binary variables represent build/no-build choices and the continuous variables represent passenger or freight flows. The quadratic integer program proposed in \cite{OKelly1987} is generally regarded as the starting point of the modern hub-location literature, and the first integer linear formulation in which hub activation, spoke allocation and inter-hub flows are treated within a unified MIP model is the one of \cite{Campbell1994Integer}. Subsequent works refined that template by exploring single-allocation versus multiple-allocation rules, by introducing tighter linearizations of the quadratic interaction terms \cite{ErnstKrishnamoorthy1996}, and by extending the framework to capacitated variants \cite{Aykin1994}. The state of the field after twenty-five years of activity is described in the comprehensive surveys \cite{AlumurKara2008,CampbellOKelly2012,FarahaniHekmatfarBolooriNikbakhsh2013}, which catalog more than a hundred formulations differing in topology assumptions, capacity rules, cost structures, and solution techniques. A common thread in those reviews is that the underlying problems become very rapidly intractable as soon as fixed costs, multicommodity flows, route feasibility, and market-level economics are represented explicitly. In airline applications this combinatorial explosion is even more apparent, because the operator's profit ultimately depends on a large number of origin-destination markets that must be served simultaneously through shared hubs and shared route capacities. To enrich the airline-specific view, structural and competitive aspects of hub-and-spoke design have been studied in \cite{AdlerSmilowitz2007}, and the operational and economic context of contemporary airline systems is summarized in \cite{BelobabaOdoniBarnhart2015}, where the strategic role of hub airports for fleet utilization, market consolidation, and revenue generation is emphasized.

\subsection{Entropy Regularization and Logit Demand Models}
A central challenge in airline network design is representing passenger demand in a behaviorally meaningful yet computationally tractable way. The simplest approach treats demand as fixed and routes it deterministically through the designed network. Such all-or-nothing assignment rules are linear and easy to embed into MIP formulations \cite{Sheffi1985,MagnantiWong1984}, but they are structurally unable to capture how real passengers react to fares, travel times, and service quality. Discrete-choice theory provides a much more meaningful description, and within that framework the multinomial logit model is by far the most widely used specification \cite{BenAkivaLerman1985,Train2009}. Formally, each passenger in market $(o,d)$ chooses an alternative $q\in\mathcal{Q}^{od}$ by maximizing a random utility
\begin{equation}
U^{od}_q=v^{od}_q+\varepsilon^{od}_q,
\label{eq:logit_random_utility}
\end{equation}
where the deterministic component $v^{od}_q$ aggregates observable disutilities such as fare and time, $v^{od}_q=\omega_p p^{od}_q+\omega_t t^{od}_q$, and the random component $\varepsilon^{od}_q$ is assumed to be independent and identically distributed Gumbel. Under those assumptions, the probability that alternative $q$ is selected in market $(o,d)$ admits the closed-form expression
\begin{equation}
P^{od}(q) = \frac{\exp\bigl(v^{od}_q\bigr)}
{\sum_{q'\in\mathcal{Q}^{od}}\exp\bigl(v^{od}_{q'}\bigr)},
\label{eq:logit_probability}
\end{equation}
which is the multinomial logit choice rule. When this expression is embedded in a network design problem so that the market shares depend on the actual network being designed, the resulting optimization model inherits nonlinear, nonconcave couplings between design variables, utilities, and demand allocation. For airline hub-and-spoke design this is precisely where realistic demand modeling starts to conflict with tractable large-scale optimization. The foundational connection between logit demand and entropy-based optimization was established by \cite{Oppenheim1993}, who showed that the logit choice probabilities arise as the unique solution of a maximum-entropy problem over the feasible share simplex. This variational viewpoint was later exploited in \cite{RealRojas2026TRB} to embed logit demand directly into a convex network design objective through utility and entropy terms, without imposing the closed-form logit equation as an explicit constraint. Existing direct remedies for the logit nonlinearity typically replace the closed-form share \eqref{eq:logit_probability} either by a piecewise-linear approximation supported on additional binary variables \cite{Vielma2015}, or by an exponential-cone reformulation \cite{LubinYamangilBentVielma2018}. Both options are valid in principle, but both also inflate the model and quickly run into the same scalability barrier that affects the underlying mixed-integer hub design problem.

\subsection{Continuous and Sparsity-Based Approximations}
The computational difficulties listed above have motivated a different modeling philosophy in which some of the binary activation decisions are replaced by sparse continuous variables. Rather than deciding directly whether an airport, hub, or service is open, one optimizes continuous capacities and then promotes sparse solutions so that inactive components collapse naturally to zero. The conceptual starting point is the well-known reweighted $\ell_1$ minimization framework, in which a nonconvex logarithmic surrogate of the $\ell_0$ pseudonorm,
\begin{equation}
\kappa\,\log\!\left(1+\frac{z}{\varepsilon}\right),
\label{eq:log_surrogate_sota}
\end{equation}
is iteratively linearized through a majorization-minimization argument, yielding at iteration $k$ the convex weighted penalty $\kappa\,z\bigl/(z^{(k-1)}+\varepsilon)$ \cite{CandesWakinBoyd2008}. Small previous values of $z^{(k-1)}$ generate large weights and aggressively penalize new activation, whereas already-active components are penalized only mildly. Iterating this update progressively drives the unused infrastructure to zero, which makes the scheme a continuous proxy of the $\ell_0$ activation cost. The same philosophy can be extended to the demand side: instead of imposing the logit expression \eqref{eq:logit_probability} as an explicit nonconvex equation, the market split can be induced through smooth entropy regularization, which preserves convexity and recovers the logit shares as optimality conditions of the resulting smooth problem. This twofold strategy—sparsity for fixed costs and entropy for demand—has been recently advocated as a scalable alternative to mixed-integer formulations for transportation network design in \cite{RealRojas2026TRB}, where it is applied to large rapid-transit network instances, and it has been extended to a single-level competitive airline hub-and-spoke design model in \cite{RealRojas2025AIAA}. Although these ideas are well established in continuous optimization and signal processing, they remain much less common in airline hub-and-spoke network design than the classical mixed-integer viewpoint.

\subsection{Bilevel Optimization and Penalty Methods}
The fourth ingredient of the proposed methodology is bilevel optimization, which is the natural mathematical framework whenever one layer of decisions is used to influence the solution of another. Following the standard convention, a bilevel program with upper-level decision $X$ and lower-level decision $y$ is written as
\begin{align}
\max_{X\in\mathcal{U}} \quad & F(X,y) \label{eq:bilevel_upper_sota}\\
\text{s.t.}\quad & y\in\arg\min_{y'\in Y(X)} f(X,y'),
\label{eq:bilevel_lower_sota}
\end{align}
where the upper-level player anticipates that the lower-level player will respond to any choice $X$ with an optimal $y$. The foundational hierarchical structure of this kind was introduced by Stackelberg \cite{vonStackelberg1934,vonStackelberg2011} in the context of competitive markets, where one firm acts as a leader by anticipating the optimal response of a follower. The first formal optimization treatment of multilevel programs within the operations research community is attributed to \cite{BrackenMcGill1973}, who showed how constraints can themselves be parametric optimization problems, and the field has been developing algorithmic and complexity theory ever since.

From the very beginning it was clear that bilevel programs are intrinsically hard. Even when both levels are linear, the feasible set of the upper-level player is in general nonconvex and disconnected, and the problem is strongly NP-hard. As a consequence, the literature has developed along several complementary directions. A first family of methods replaces the lower-level optimization by its Karush-Kuhn-Tucker (KKT) conditions, producing a single-level mathematical program with complementarity constraints whose binary indicators select the active inequalities. This route works particularly well when the lower-level problem is convex, and it has been widely used since the early references on quadratic bilevel programming. A second family adapts descent methods to the bilevel structure by approximating the hyperderivative through the implicit-function theorem; in recent years, these implicit-gradient descent (IGD) ideas have become the backbone of bilevel machine-learning algorithms \cite{LuMei2024,Jiang2024Primal,GhadimiWang2018}. A third family relies on evolutionary or metaheuristic strategies and is intended for problems whose mathematical structure is too unfavorable for derivative-based methods \cite{SinhaMaloDeb2018}. Standard survey references covering these developments include \cite{Bard1998,ColsonMarcotteSavard2007,Dempe2002,SinhaMaloDeb2018}.

Penalty-based reformulations occupy a central place in this landscape and are particularly relevant for the present work. The basic idea, due to Aiyoshi and Shimizu, is to replace the lower-level optimality condition by a penalty term added to the upper-level objective \cite{AiyoshiShimizu1981,AiyoshiShimizu1984}. When the lower level is unconstrained, a typical penalty is $\|\nabla_y g(u,y)\|^2$, so that large penalty weights drive the gradient of the lower-level objective to zero. When the lower level has constraints, the same idea is applied to the gradient of the associated Lagrangian together with the remaining KKT conditions. An alternative formulation that is more convenient for the algorithms used in this work replaces the lower-level optimality by a \emph{value-function} penalty, namely $g(u,y)-v(u)$, where $v(u)$ denotes the lower-level optimal value. Exact penalization results for this construction were established in \cite{Ye1997Exact}, and recent first-order penalty algorithms have shown that constrained bilevel problems with convex lower-level structure can be approached through penalized single-level surrogates with implementable complexity guarantees \cite{LuMei2024,Jiang2024Primal}. The algorithmic scheme developed in the present paper belongs to that family.

To close, it is worth stressing that the present paper lies at the intersection of all four strands above. Relative to the traditional hub-location literature, it avoids an explicit mixed-integer treatment of airport, hub, and service activation and instead works with a continuous sparse design model. Relative to standard deterministic network-design formulations, it incorporates endogenous demand capture through utility and entropy terms, so that market shares react to the designed network. Relative to generic sparsity-based continuous models, it is tailored to airline hub-and-spoke design and to a profit-oriented private operator. Finally, relative to the general bilevel literature, the upper level is not used here to redesign the network directly, but to calibrate the market-specific coefficients that guide the lower-level continuous approximation toward the airline's true profit objective. This combination of entropy-based demand modeling, sparse continuous activation, and value-function-penalty bilevel calibration is, to our knowledge, not standard in the airline hub-and-spoke design literature and defines the methodological contribution developed in the remainder of the paper.
